% ============================================================
%  THE PENULTIMATE STATE:
%  Why the Reversed Configuration Is Not the Initial
%  Configuration, and a Caesium-Ion Validation
% ============================================================

\documentclass[twocolumn,10pt,a4paper]{article}

\usepackage[utf8]{inputenc}
\usepackage[T1]{fontenc}
\usepackage{lmodern}
\usepackage[a4paper,top=2cm,bottom=2.2cm,left=1.8cm,right=1.8cm,columnsep=0.6cm]{geometry}
\usepackage{amsmath,amssymb,amsthm,mathtools}
\usepackage{bm}
\usepackage{booktabs}
\usepackage{array}
\usepackage{enumitem}
\usepackage{graphicx}
\usepackage{xcolor}
\usepackage{microtype}
\usepackage[round,sort&compress]{natbib}
\usepackage[colorlinks=true,linkcolor=blue!60!black,
            citecolor=blue!60!black,urlcolor=blue!60!black]{hyperref}
\usepackage{fancyhdr}
\usepackage{algorithm}
\usepackage{algpseudocode}
\usepackage{float}

\pagestyle{fancy}
\fancyhf{}
\fancyhead[LE,RO]{\thepage}
\fancyhead[RE]{\small The Penultimate State}
\fancyhead[LO]{\small Sachikonye}
\renewcommand{\headrulewidth}{0.4pt}

\newtheoremstyle{thmstyle}{\topsep}{\topsep}{\itshape}{}{\bfseries}{.}{.5em}{}
\theoremstyle{thmstyle}
\newtheorem{theorem}{Theorem}[section]
\newtheorem{lemma}[theorem]{Lemma}
\newtheorem{proposition}[theorem]{Proposition}
\newtheorem{corollary}[theorem]{Corollary}
\theoremstyle{definition}
\newtheorem{definition}[theorem]{Definition}
\newtheorem{axiom}{Axiom}
\newtheorem{protocol}[theorem]{Protocol}
\theoremstyle{remark}
\newtheorem{remark}[theorem]{Remark}
\newtheorem{principle}{Principle}

\newcommand{\R}{\mathbb{R}}
\newcommand{\N}{\mathbb{N}}
\newcommand{\Zp}{\mathbb{Z}^{+}}
\newcommand{\Parts}{\mathcal{P}}
\newcommand{\PhiMap}{\varPhi}
\newcommand{\Sspace}{\mathcal{S}}
\newcommand{\Depth}{\mathcal{M}}
\newcommand{\Lagr}{\mathcal{L}}
\newcommand{\partinert}{\mu}
\newcommand{\Apartition}{\mathbf{A}_{\Depth}}
\newcommand{\kB}{k_{\mathrm{B}}}
\newcommand{\calF}{\mathcal{F}}

\title{\textbf{The Penultimate State as Termination Endpoint of Unconstrained Sub-task Recursion :Operational Distinctions of State Despite Configuration Reversal Distinguishable }}

\author{
Kundai Farai Sachikonye\\
Technical University of Munich\\
Chair of Experimental Bioinformatics\\
AIMe Registry for Artificial Intelligence\\
\texttt{kundai.sachikonye@wzw.tum.de}
}

\date{\today}

\begin{document}
\maketitle

\begin{abstract}
\noindent
We argue that Loschmidt's paradox dissolves under a single observation:
a configuration is a structure of correlations, and structures of
correlations accumulate monotonically under any traversal direction.
Velocity reversal moves a system through configurations that resemble
earlier ones, but the cumulative configuration count is independent of
velocity sign. We develop this thesis from first principles, deriving
every required result in self-contained form: a bounded-phase-space
axiom; oscillatory necessity; the partition coordinate system
$(n,\ell,m,s)$ with capacity $C(n)=2n^{2}$; a bijection
$\PhiMap:\Zp\to\Parts$ between the positive integers and the
configuration states; the time-count identity $dM/dt=\omega/(2\pi)$ as
a definition of derived time; the partition Lagrangian whose
specialisations yield all four canonical analyser equations of mass
spectrometry; the snapshot/velocity decoupling that makes the
cumulative count monotone in any direction; and a backward
trajectory-completion algorithm whose termination at the
\emph{penultimate} state is the structural requirement for
recognition. The argument exhibits two self-undermining structures
in the conventional reading of Loschmidt's reversal. First, the SI
second is defined as $9{,}192{,}631{,}770$ velocity reversals of
caesium-133 per stipulated unit; if reversal cancelled rather than
accumulated, the unit could not exist. Second, the same atom, used as
an analyte in an Orbitrap, performs roughly $10^{6}$ velocity reversals
per second of analytical observation; the partition Lagrangian
predicts the partition count $M(T)=\omega_{z}T/(2\pi)$ to be a strictly
monotone function of $T$, and the cumulative phase recovered from the
Orbitrap image current via Hilbert transform measures it to machine
precision. We specify this experiment in full operational detail.
Supporting predictions include: a bijection check on standard
\texttt{mzML} ions; verification of the four analyser scaling laws
under the single Lagrangian; the subharmonic encoding of MS/MS fragment
spectra in MS\textsuperscript{1} transients; and a backward-completion
complexity check showing $\log_{3}N$ steps under virtual sub-state
decomposition and $\Theta(N)$ collapse without. No prior work of the
author is cited; the literature references are restricted to standard
mathematical and physical lemmas used in the derivations.

\medskip
\noindent\textbf{Keywords:} Loschmidt's paradox, configuration
monotonicity, partition coordinates, time-count identity, snapshot
decoupling, caesium definition of the second, Orbitrap transient,
penultimate state, backward trajectory completion, falsifiable
mass-spectrometry validation.
\end{abstract}

\tableofcontents
\bigskip\hrule\bigskip

% =========================================================
\section{Introduction}
\label{sec:intro}
% =========================================================

\subsection{The single claim}

This paper makes one claim and supports it operationally.

\begin{quote}
\emph{The reversed configuration is not the initial configuration.
The configuration count is monotone under any traversal direction.}
\end{quote}

Equivalently: Loschmidt's paradox \citep{loschmidt1876,boltzmann1877}
identifies as ``the same state'' two physical situations that are
not the same. They share a projection onto $(q,p)$ coordinates, but
they differ in a separate physical coordinate---the cumulative
configuration count---that the conventional reading does not
acknowledge as a coordinate at all.

Once the coordinate is acknowledged, the paradox dissolves. The work
of this paper is (i) to derive every result needed to make the
acknowledgement rigorous, (ii) to show that the cumulative
configuration count is well defined and monotone under any traversal
direction, and (iii) to validate the framework operationally on
a caesium ion in an Orbitrap mass spectrometer, where roughly
$10^{6}$ velocity reversals occur during a single transient and the
configuration count is predicted to advance monotonically through
every one of them.

\subsection{Two self-undermining structures in Loschmidt's reading}

The conventional reading of Loschmidt's velocity reversal makes
predictions that are inconsistent with two pieces of standard
physics the reading itself accepts.

\textbf{The metrology argument.} The SI second is defined as the
duration of $9{,}192{,}631{,}770$ periods of the radiation
corresponding to the transition between the two hyperfine levels of
the ground state of caesium-133. Each period contains two velocity
reversals of the relevant electronic motion: the electron moves one
way, decelerates, reverses, moves the other way, decelerates, reverses
again. The international definition of time is built on the empirical
fact that these reversals \emph{accumulate} rather than cancel: every
reversal advances the count by exactly one half-cycle, and the unit
``second'' is the named duration of $9{,}192{,}631{,}770$ such
advancements. If reversal cancelled count, this number would equal
zero, and the SI second would be undefined. The unit exists. Therefore
reversal accumulates.

\textbf{The decoupling argument.} A configuration of a bounded
oscillatory system is a topological structure: a partition cell, a
phase-lock network, a categorical address. It is defined by the
arrangement of correlations among constituents. Velocity is the rate
at which the system moves between configurations; it is not the cause
of any particular configuration. The same configuration can be reached
at any velocity, including its reverse. Therefore reversing velocities
does not undo configurations; it traverses configurations resembling
earlier ones, and each such traversal contributes to the cumulative
count.

These two arguments are independent but mutually reinforcing. The
metrology argument is operational: the conventional reading of
Loschmidt would null the unit in which the reading is stated. The
decoupling argument is structural: configurations and velocities are
different objects, and the cumulative count is a property of the
former.

\subsection{Backward trajectory completion and the recognition gap}

The framework's positive content is summarised by an algorithm:
\emph{backward trajectory completion}. Given a known endpoint
configuration at depth $k$ in a ternary refinement hierarchy of size
$N=3^{k}$, the algorithm navigates from endpoint to root in exactly
$\log_{3}N$ steps. Critically, the algorithm terminates at the
\emph{penultimate} state---one isomorphism step short of the endpoint
itself---because that gap is the structural requirement for
recognition. To complete the trajectory all the way to the endpoint
would be to become the endpoint, and recognition requires a
distinction between the recognising state and the recognised state.

Loschmidt's reversal is the question of whether this gap can be
closed. It cannot, for the same reason: closing it would dissolve
the recognition the operation presupposes. The framework's algorithm
exhibits the gap as a constitutive feature; the operational experiment
on Cs\textsuperscript{+} in an Orbitrap exhibits the cumulative
count's monotonicity even as roughly $10^{6}$ velocity reversals
occur during a single transient.

\subsection{Plan of the paper}

Section~\ref{sec:axiom} establishes the single axiom and derives
oscillatory necessity. Section~\ref{sec:partition} derives the
partition coordinate system and the capacity formula.
Section~\ref{sec:bijection} proves the bijection $\PhiMap:\Zp\to\Parts$.
Section~\ref{sec:timecount} derives the time-count identity as a
non-circular definition of derived time.
Section~\ref{sec:configuration} states the configuration-as-snapshot
argument and proves the decoupling of configuration from velocity.
Section~\ref{sec:lagrangian} derives the partition Lagrangian and the
four analyser scaling laws.
Section~\ref{sec:dissolution} dissolves Loschmidt's paradox: the
reversed configuration is not the initial configuration; the
cumulative count is monotone; the SI second is the metrological lock.
Section~\ref{sec:occspec} diagnoses why the conventional reading
nevertheless feels paradoxical: it conflates the cost an embedded
agent pays to \emph{specify} a low-probability state with an
impossibility of the system's \emph{occupying} one, and it commits
three category errors that the occupation--specification asymmetry
blocks.
Section~\ref{sec:backward} introduces the backward-completion algorithm
and the penultimate-state termination, deriving the recognition gap
from the resolution-floor theorem rather than from informal grounds.
Section~\ref{sec:phase} proves the phase-coherence theorem and the
subharmonic-encoding consequence for MS/MS spectra.
Section~\ref{sec:experiment} specifies the headline experiment:
Cs\textsuperscript{+} in an Orbitrap with Hilbert-transform phase
recovery.
Section~\ref{sec:supporting} states four supporting experiments.
Section~\ref{sec:discussion} discusses limitations and scope.

% =========================================================
\section{The Single Axiom and Oscillatory Necessity}
\label{sec:axiom}
% =========================================================

\subsection{Bounded phase space}

Consider a classical dynamical system with $d$ spatial degrees of
freedom on a phase space $\Gamma\subseteq\R^{2d}$ with canonical
coordinates $(\mathbf{q},\mathbf{p})$. A Hamiltonian flow
$\phi_{t}:\Gamma\to\Gamma$ preserves the Liouville measure
$\mu=d^{d}q\,d^{d}p/h^{d}$ for a fixed phase-space cell size $h$
\citep{arnold1989}.

\begin{axiom}[Bounded Phase Space]
\label{ax:bps}
Every persistent dynamical system occupies a bounded region
$\Omega\subseteq\Gamma$ of finite Liouville measure
$\mu(\Omega)<\infty$, and $\Omega$ admits hierarchical partitioning
into distinguishable subregions.
\end{axiom}

A \emph{persistent} system is one whose categorical identity is
preserved over time. The motivation is empirical: an electron in an
atomic well, a vibrational mode of a molecule, an ion in an
electromagnetic trap, a photon in a cavity. Each is confined; the
unbounded idealisation does not persist.

\subsection{Poincar\'{e} recurrence}

\begin{theorem}[Poincar\'{e} Recurrence]
\label{thm:poincare}
Let $(\Omega,\mu,\phi_{t})$ be measure-preserving with
$\mu(\Omega)<\infty$. For measurable $A\subseteq\Omega$ with
$\mu(A)>0$, almost every $x\in A$ returns to $A$ infinitely often.
\end{theorem}

\begin{proof}
Fix $T>0$ and set $B=\{x\in A:\phi_{nT}(x)\notin A\ \forall n\geq 1\}$.
The sets $\{\phi_{nT}(B)\}_{n\geq 0}$ are pairwise disjoint: if
$\phi_{jT}(B)\cap\phi_{kT}(B)\neq\emptyset$ for $j<k$ then some
$x\in B$ would have $\phi_{(k-j)T}(x)\in A$, contradicting the
definition of $B$. By measure preservation,
$\mu(\phi_{nT}(B))=\mu(B)$. If $\mu(B)>0$ then
$\sum_{n\geq 0}\mu(\phi_{nT}(B))=\infty>\mu(\Omega)$, a contradiction.
Hence $\mu(B)=0$.
\end{proof}

\subsection{Oscillatory necessity}

\begin{theorem}[Oscillatory Necessity]
\label{thm:osc}
The only class of dynamics on a bounded measure-finite phase space
supporting categorical observation is oscillatory.
\end{theorem}

\begin{proof}
We eliminate three alternatives.

\emph{Static}: if $\phi_{t}(x)=x$ for all $t$ and $x$, no temporal
structure is present and no measurement can occur.

\emph{Monotonic}: if there exists $f:\Omega\to\R$ with
$t\mapsto f(\phi_{t}(x))$ strictly monotonic, then for any
$\varepsilon>0$ the set $A_{\varepsilon}=\{x:f(x)<\varepsilon\}$ is
exited once and not revisited, contradicting
Theorem~\ref{thm:poincare} unless $\mu(A_{\varepsilon})=0$ for all
$\varepsilon$, contradicting $\mu(\Omega)<\infty$.

\emph{Chaotic}: if
$|\phi_{t}(x)-\phi_{t}(y)|\sim|x-y|e^{\lambda t}$ with $\lambda>0$,
the categorical labels are not stable; the same preparation produces
different labels in finite time, violating reproducibility.

Oscillatory motion returns to neighbourhoods of previous states by
Theorem~\ref{thm:poincare}; it is bounded, recurrent, and
identity-preserving.
\end{proof}

\begin{corollary}[Discrete Mode Spectrum]
\label{cor:modes}
Bounded oscillatory dynamics admits a discrete frequency spectrum,
the eigenvalues of the Koopman operator $U_{t}:f\mapsto f\circ\phi_{t}$
on $L^{2}(\Omega,\mu)$ \citep{koopman1931}.
\end{corollary}

\begin{proof}
$U_{t}$ is unitary on $L^{2}(\Omega,\mu)$; on a finite-measure compact
phase space it has compact resolvent and therefore discrete spectrum.
\end{proof}

\begin{corollary}[Frequency--Energy Identity]
\label{cor:Ehbar}
A fundamental oscillatory mode of angular frequency $\omega$ in a
bounded domain satisfies $E=\hbar\omega$.
\end{corollary}

\begin{proof}
Bohr--Sommerfeld quantisation
$\oint p\,dq=n\hbar\cdot 2\pi$ from single-valuedness on a bounded
domain \citep{bohr1913,landau1976} gives $E_{n}=n\hbar\omega$. The
fundamental mode is $n=1$.
\end{proof}

% =========================================================
\section{Partition Coordinates and Capacity}
\label{sec:partition}
% =========================================================

\subsection{The four coordinates}

\begin{theorem}[Partition Coordinate System]
\label{thm:nlms}
A hierarchically nested three-dimensional bounded oscillatory system
admits a complete coordinate system $(n,\ell,m,s)$ with
\begin{align*}
n&\in\Zp,&\ell&\in\{0,1,\ldots,n-1\},\\
m&\in\{-\ell,\ldots,+\ell\},& s&\in\{-\tfrac{1}{2},+\tfrac{1}{2}\}.
\end{align*}
\end{theorem}

\begin{proof}
\emph{$n$:} nested shells $\Omega_{1}\supset\Omega_{2}\supset\cdots$
provided by Axiom~\ref{ax:bps} label the principal depth.

\emph{$\ell$:} at depth $n$ the oscillatory boundary supports angular
modes with wavevector $k_{\theta}=\ell/R_{n}$ for shell radius
$R_{n}\propto n$; the mode fits the radial extent only if
$\ell\leq n-1$.

\emph{$m$:} the orientation of the angular mode in three-space is
labelled by an irreducible representation of $SO(3)$; the irreducible
representation of angular momentum $\ell$ has dimension $2\ell+1$,
basis $m\in\{-\ell,\ldots,+\ell\}$ \citep{hall2015}.

\emph{$s$:} the fundamental group $\pi_{1}(SO(3))=\mathbb{Z}_{2}$
\citep{nakahara2003} gives two topologically distinct boundary
helicities, labelled $s=\pm\tfrac{1}{2}$.
\end{proof}

\subsection{The capacity formula}

\begin{theorem}[Shell Capacity]
\label{thm:cap}
The number of distinct partition states at depth $n$ is
$C(n)=2n^{2}$.
\end{theorem}

\begin{proof}
$C(n)=\sum_{\ell=0}^{n-1}(2\ell+1)\cdot 2=2\sum_{\ell=0}^{n-1}(2\ell+1)
=2n^{2}$, using $\sum_{\ell=0}^{n-1}(2\ell+1)=n^{2}$.
\end{proof}

\begin{corollary}[Cumulative State Count]
\label{cor:cum}
The total number of states with principal depth $\leq N$ is
\begin{equation}
N_{\mathrm{state}}(N)=\sum_{n=1}^{N}2n^{2}=\frac{N(N+1)(2N+1)}{3}.
\label{eq:Ns}
\end{equation}
\end{corollary}

\begin{proof}
Standard sum-of-squares identity multiplied by two.
\end{proof}

% =========================================================
\section{The Partition Bijection}
\label{sec:bijection}
% =========================================================

\subsection{The map}

Let $\Parts$ be the set of partition states with coordinates in the
ranges of Theorem~\ref{thm:nlms}.

\begin{definition}[Partition Map]
\label{def:phi}
$\PhiMap:\Zp\to\Parts$ is defined by: given $M\in\Zp$,
\begin{enumerate}[nosep,leftmargin=*]
\item find $n$ with $N_{\mathrm{state}}(n-1)<M\leq N_{\mathrm{state}}(n)$;
\item set $\delta=M-N_{\mathrm{state}}(n-1)\in\{1,\ldots,2n^{2}\}$;
\item enumerate level-$n$ states lexicographically and assign the
$\delta$-th tuple as $(\ell(M),m(M),s(M))$;
\item set $\PhiMap(M)=(n(M),\ell(M),m(M),s(M))$.
\end{enumerate}
\end{definition}

\begin{theorem}[Bijection]
\label{thm:bij}
$\PhiMap$ is a bijection.
\end{theorem}

\begin{proof}
\emph{Injectivity:} $\PhiMap(M_{1})=\PhiMap(M_{2})$ forces
$n(M_{1})=n(M_{2})$ and $\delta_{1}=\delta_{2}$ by the strict
monotonicity of the lexicographic enumeration; hence $M_{1}=M_{2}$.

\emph{Surjectivity:} for $(n,\ell,m,s)\in\Parts$, set
\[
M_{0}=N_{\mathrm{state}}(n-1)+2\sum_{\ell'=0}^{\ell-1}(2\ell'+1)
+2(m+\ell)+(s+\tfrac{1}{2}).
\]
Then $N_{\mathrm{state}}(n-1)<M_{0}\leq N_{\mathrm{state}}(n)$ and the
level-$n$ enumeration assigns $(n,\ell,m,s)$ to $M_{0}$.
\end{proof}

\begin{corollary}[Invertibility]
\label{cor:inv}
$\PhiMap^{-1}$ is the explicit formula for $M_{0}$ above.
\end{corollary}

\subsection{Traversal versus existence}

\begin{definition}[Forward Sequence]
\label{def:forward}
$\calF(M_{t})=\{\PhiMap(M_{t}+k):k\geq 1\}\subseteq\Parts$.
\end{definition}

\begin{theorem}[Successor Existence]
\label{thm:succ}
For every $M_{t}\in\Zp$, the partition state $\PhiMap(M_{t}+1)$ is a
well-defined member of $\Parts$, independently of whether the system
has traversed $M_{t}+1$.
\end{theorem}

\begin{proof}
$M_{t}+1\in\Zp$ by Peano \citep{peano1889}; $\PhiMap$ is defined on
all of $\Zp$ by Theorem~\ref{thm:bij}.
\end{proof}

\begin{figure*}[!htbp]
    \centering
    \includegraphics[width=\textwidth]{figures/panel_02_bijection.png}
    \caption{\textbf{Partition Coordinates, Capacity, and the Integer-Indexed Bijection.}
    (\textbf{A})~Shell capacity $C(n)$ on the linear $y$-axis ($0$ to $1800$) versus principal depth $n$ on the linear $x$-axis ($1$ to $30$), enumerated directly over $(\ell,m,s)$ tuples with $\ell\in\{0,\ldots,n-1\}$, $m\in\{-\ell,\ldots,+\ell\}$, $s\in\{-1/2,+1/2\}$. The enumeration agrees with $C(n)=2n^{2}$ at every $n$; maximum residual zero. Empirical content of Theorem~\ref{thm:cap} (Shell Capacity).
    (\textbf{B})~Cumulative state count on a log $y$-axis ($10^{0}$ to $10^{5}$) versus depth bound $N$ on the linear $x$-axis ($1$ to $30$). Steelblue circles: $N_{\mathrm{state}}(N)$ obtained by summing the enumerated capacities. Crimson curve: the closed-form prediction $N(N{+}1)(2N{+}1)/3$. The two coincide at every $N$ with zero residual. Empirical content of Corollary~\ref{cor:cum} (Cumulative State Count) and the foundation for the bijection in panels~C and~D.
    (\textbf{C})~Selection rules visualised: $m$ on the linear $y$-axis ($-12$ to $+12$) versus $\ell$ on the linear $x$-axis ($0$ to $12$), $2000$ integer indices $M\in[1,2000]$ each forward-mapped through $\PhiMap$. Marker colour encodes the spin coordinate $s\in\{-1/2,+1/2\}$ via a diverging coolwarm scale. The point cloud occupies exactly the triangular region $|m|\leq\ell$, with no excursion above $|m|=\ell$ and no missing $(\ell,m)$ pair for any reachable $\ell$. Zero capacity violations and zero selection violations across $2000$ samples. Empirical content of Theorem~\ref{thm:bij} (Bijection).
    (\textbf{D})~Three-dimensional scatter of $(n,\ell,m)$, viewing angle elevation $30^{\circ}$ azimuth $-60^{\circ}$, $2000$ points from the same $M$-range coloured by $s$ on the coolwarm scale. Each occupied $(n,\ell,m,s)$ tuple is exhibited as a distinct cluster; no point appears twice for distinct $M$. Combined with the $220\,000$-sample round-trip test of E2 (forward failures: $0$; backward failures: $0$; capacity violations: $0$; selection violations: $0$), this is the empirical realisation of $\PhiMap:\Zp\to\Parts$ as an honest bijection.}
    \label{fig:panel-bijection}
\end{figure*}

% =========================================================
\section{Time as Configuration Count}
\label{sec:timecount}
% =========================================================

\subsection{Closure as a topological primitive}

\begin{definition}[Oscillatory Closure]
\label{def:closure}
An \emph{oscillatory closure} is the event of an angular coordinate
$\theta$ completing one traversal of $[0,2\pi)$ and returning to its
initial value.
\end{definition}

\begin{proposition}[Closure is Dimensionless]
\label{prop:closure}
A closure is definable without reference to duration. It requires (i)
a configuration space carrying an angular coordinate $\theta$, and
(ii) the predicate that $\theta$ has traversed $[0,2\pi)$ once.
Neither invokes a temporal unit.
\end{proposition}

\begin{proof}
$[0,2\pi)$ is dimensionless angular extent. The predicate of return
to initial value is a property of trajectory geometry, not of duration.
Poincar\'{e} recurrence (Theorem~\ref{thm:poincare}) guarantees that
such returns occur on bounded oscillatory systems.
\end{proof}

\subsection{Time as derived parameter}

\begin{definition}[Time as Count]
\label{def:tcount}
For a bounded oscillatory system with an attached reference oscillator
at frequency $f_{\mathrm{ref}}$ stipulated by convention, the
\emph{time} of a partition count $M$ is the derived ratio
\begin{equation}
t\;:=\;\frac{M}{f_{\mathrm{ref}}}.
\label{eq:tratio}
\end{equation}
\end{definition}

\begin{proposition}[Non-Circularity]
\label{prop:nocirc}
Equation \eqref{eq:tratio} is not circular. $M$ counts topological
closures (Proposition~\ref{prop:closure}), each dimensionless;
$f_{\mathrm{ref}}$ is a stipulated integer (e.g., $f_{\mathrm{Cs}}=
9{,}192{,}631{,}770$ defining the SI second \citep{si2019}). The
right-hand side is a ratio of pure integers, named $t$ in the
stipulated unit. No prior notion of duration is required.
\end{proposition}

\begin{theorem}[Time-Count Identity]
\label{thm:tc}
For a bounded oscillatory system of angular frequency $\omega$,
\begin{equation}
\boxed{\;\frac{dM}{dt}=\frac{\omega}{2\pi}=\frac{1}{\tau_{p}},\;}
\label{eq:tc}
\end{equation}
where $\tau_{p}=2\pi/\omega$ is the residence time per state.
\end{theorem}

\begin{proof}
One cycle traverses one state per period $T_{p}=2\pi/\omega$.
\end{proof}

\begin{corollary}[Strict Monotonicity]
\label{cor:mono}
$M(t_{2})>M(t_{1})$ for $t_{2}>t_{1}$.
\end{corollary}

\begin{corollary}[Exact Phase]
\label{cor:phase}
The oscillation phase at count $M$ is $\theta(M)=2\pi M$.
\end{corollary}

\begin{proof}
$\theta=\omega t=2\pi f_{\mathrm{ref}}\cdot M/f_{\mathrm{ref}}=2\pi M$
for a reference oscillator at integer frequency.
\end{proof}

\subsection{The SI second as metrological lock}

The definition of the SI second is the empirical statement that
velocity reversals of the caesium-133 hyperfine transition accumulate
rather than cancel. The unit ``second'' is the named duration of
$9{,}192{,}631{,}770$ such cycles \citep{si2019}, each containing two
velocity reversals of the relevant electronic motion. The unit
exists, hence the count accumulates. This is not a derivation; it is
an empirical observation incorporated into international metrology.
We return to it as the operational lock in Section~\ref{sec:dissolution}.

% =========================================================
\section{Configurations are Snapshots; Velocities are Rates}
\label{sec:configuration}
% =========================================================

\subsection{Configurations}

\begin{definition}[Configuration]
\label{def:config}
A \emph{configuration} of a bounded $N$-body system at time $t$ is
the structural specification of its inter-constituent correlations:
the equivalence class of phase-space points sharing the same
partition state $(n,\ell,m,s)$ together with the same network of
inter-particle correlations.
\end{definition}

A configuration is a structure, not a kinematic state. It records
\emph{which constituents are correlated with which others, in what
patterns}. It does not record \emph{how fast} the constituents are
moving.

\begin{definition}[Configuration Count]
\label{def:Cc}
The \emph{configuration count} $C(t)$ at time $t$ is the cumulative
number of distinct configurations the system has occupied over its
history up to $t$, indexed by partition count $M$ via
Theorem~\ref{thm:bij}.
\end{definition}

By Corollary~\ref{cor:mono}, $C(t)$ is monotone non-decreasing in $t$.

\subsection{Decoupling of configuration from velocity}

\begin{proposition}[Same Configuration at Different Velocities]
\label{prop:decouple}
A given configuration can be occupied by a system at any velocity
distribution consistent with the bounded phase space $\Omega$. The
configuration $(n,\ell,m,s)$ specifies the structural arrangement of
correlations; it does not specify the magnitude or direction of any
constituent's velocity.
\end{proposition}

\begin{proof}
The coordinates $(n,\ell,m,s)$ are topological labels of the
oscillatory boundary structure (Theorem~\ref{thm:nlms}). The boundary
admits the same configuration at any speed of traversal: the same
shell, same angular nodal structure, same orientation, same chirality
can be occupied while the constituents move slowly or quickly. The
configuration is determined by where the boundary \emph{is}, not by
how fast anything is moving through it.
\end{proof}

\begin{remark}
\label{rem:tempindep}
Proposition~\ref{prop:decouple} entails that the same configuration
can be present at arbitrary kinetic energy distributions, including
the same configuration at different temperatures. Temperature is a
property of the velocity distribution at a configuration; the
configuration is not a property of temperature. This is the decoupling
that makes the cumulative count monotone under any traversal direction.
\end{remark}

\subsection{Monotonicity under arbitrary traversal direction}

\begin{theorem}[Configuration-Count Monotonicity]
\label{thm:mono}
Let a bounded oscillatory system evolve under any sequence of
velocity-direction choices, including reversals applied at arbitrary
moments. Then the configuration count $C(t)$ is strictly monotone
non-decreasing in $t$.
\end{theorem}

\begin{proof}
By Proposition~\ref{prop:decouple}, the configuration occupied by the
system at any moment is determined by the structural state of the
oscillatory boundary, not by the velocity. By
Theorem~\ref{thm:poincare}, bounded dynamics is recurrent: the system
will occupy configurations resembling earlier ones. By
Theorem~\ref{thm:bij}, each occupation event corresponds to a distinct
$M\in\Zp$, distinct from all previous $M$. The map
$M\mapsto\PhiMap(M)$ is many-to-one onto the structural type
$(n,\ell,m,s)$: two distinct integers $M_{1}\neq M_{2}$ can yield
configurations of the same structural type at different times.
Therefore $C(t)$, which counts \emph{occupation events} indexed by
$M\in\Zp$, increases by one at each event, irrespective of whether
the structural type at the event resembles a structural type at an
earlier event.

Velocity reversal applied at $\tau$ causes the system to move through
configurations whose structural types resemble those at $t<\tau$, but
each such occupation is a distinct event in $\Zp$ and increments
$C$. The reversed traversal does not subtract; it adds.
\end{proof}

\begin{corollary}[The Reversed State is Not the Initial State]
\label{cor:notinitial}
Let a system at time $0$ be in configuration $\PhiMap(M_{0})$. Suppose
velocities are reversed at time $\tau$ and the system evolves until
time $2\tau$ to a configuration whose structural type
$(n,\ell,m,s)$ matches that at time $0$. Then the partition count at
time $2\tau$ satisfies $M_{2\tau}=M_{0}+2N_{\tau}>M_{0}$, where
$N_{\tau}$ is the number of configurations traversed during $[0,\tau]$.
\end{corollary}

\begin{proof}
Directly from Theorem~\ref{thm:mono} applied over $[0,2\tau]$.
\end{proof}

\begin{figure*}[!htbp]
    \centering
    \includegraphics[width=\textwidth]{figures/panel_01_monotonicity.png}
    \caption{\textbf{Configuration-Count Monotonicity Under Arbitrary Velocity Reversals.}
    (\textbf{A})~Position $x(t)$ on the linear $y$-axis ($-1$ to $+1$) versus time on the linear $x$-axis ($0$ to $4$~s) for a bounded harmonic oscillator at $\omega/(2\pi)=5$~Hz, sampled at $f_{s}=2000$~Hz. Forty velocity reversals are injected at uniformly random times inside the central seven-eighths of the window (faint crimson vertical markers); the trajectory remains bounded and the recurrent topology is preserved through every reversal. This is the kinematic substrate of the experiment: the $(q,p)$ projection visibly retraces near each reversal.
    (\textbf{B})~Cumulative partition count on the linear $y$-axis ($0$ to $20$) versus time on the linear $x$-axis ($0$ to $4$~s). Steelblue: the framework count $M(t)$ from the time-count identity~\eqref{eq:tc}, advancing by $\omega\,dt/(2\pi)$ at every step irrespective of velocity direction. Crimson: the conventional Loschmidt prediction in which velocity reversal subtracts from the count. The two curves coincide until the first reversal and then diverge; over the full window the framework attains $M(T)=20$ while the conventional curve oscillates near $6$, never recovering. Empirical content of Theorem~\ref{thm:mono} (Configuration-Count Monotonicity).
    (\textbf{C})~Turning-point detail: cumulative count on the linear $y$-axis versus time on the linear $x$-axis in a $0.2$~s window centred on a single reversal (neutral dashed vertical reference). The framework curve passes through the reversal with no inflection, no decrement, and no perturbation of slope $\omega/(2\pi)$; the conventional curve kinks downward. This is the empirical content of Corollary~\ref{cor:notinitial} (The Reversed State is Not the Initial State): the reversal does not subtract from the cumulative count, even at the precise instant of reversal.
    (\textbf{D})~Three-dimensional scatter of $(q,p,M)$ over the entire trajectory, viridis colourmap by $t$, viewing angle elevation $30^{\circ}$ azimuth $-60^{\circ}$. The $(q,p)$ projection lies in the unit disc and is revisited many times; the $M$ coordinate increases strictly along the orbit, lifting each revisit to a distinct height. The trajectory traces a spiral whose projection on $(q,p)$ is recurrent and whose lift in $M$ is monotone. Empirical content of Theorem~\ref{thm:mono} stated geometrically: the cumulative count is the missing coordinate the conventional reading does not acknowledge.}
    \label{fig:panel-monotonicity}
\end{figure*}

% =========================================================
\section{The Partition Lagrangian}
\label{sec:lagrangian}
% =========================================================

\subsection{Construction}

\begin{theorem}[Partition Lagrangian]
\label{thm:Lag}
A bounded oscillatory mode with inertia $\partinert=\alpha(m/z)$ in a
scalar depth landscape $\Depth(\mathbf{x},t)$ and vector potential
$\Apartition(\mathbf{x})$ admits the unique Galilean-invariant
Lagrangian
\begin{equation}
\boxed{\;\Lagr_{\Depth}=\tfrac{1}{2}\partinert|\dot{\mathbf{x}}|^{2}
+\partinert\dot{\mathbf{x}}\cdot\Apartition-\Depth.\;}
\label{eq:Lagdef}
\end{equation}
\end{theorem}

\begin{proof}
Galilean invariance forces a quadratic-in-velocity kinetic term, the
unique scalar being $\tfrac{1}{2}\partinert|\dot{\mathbf{x}}|^{2}$.
Minimal couplings to scalar and vector fields are $-\Depth$ and
$\partinert\dot{\mathbf{x}}\cdot\Apartition$ respectively. Higher
derivatives are excluded by the Ostrogradsky instability
\citep{ostrogradsky1850}.
\end{proof}

\begin{theorem}[Equations of Motion]
\label{thm:EOM}
The Euler--Lagrange equations of \eqref{eq:Lagdef} are
\begin{equation}
\partinert\ddot{\mathbf{x}}=-\nabla\Depth
+\partinert\dot{\mathbf{x}}\times\mathbf{B}_{\Depth},
\quad
\mathbf{B}_{\Depth}=\nabla\times\Apartition.
\label{eq:EL}
\end{equation}
\end{theorem}

\begin{proof}
Standard Euler--Lagrange algebra on \eqref{eq:Lagdef}, with the
magnetic-like term obtained from the vector identity
$\nabla(\dot{\mathbf{x}}\cdot\Apartition)-d\Apartition/dt
=\dot{\mathbf{x}}\times(\nabla\times\Apartition)$.
\end{proof}

\begin{corollary}[Lorentz Structure]
\label{cor:lor}
Setting $\Depth=q\phi$ and $\Apartition=(q/\partinert)\mathbf{A}$
reduces \eqref{eq:EL} to
$\partinert\ddot{\mathbf{x}}=q\mathbf{E}+q\dot{\mathbf{x}}\times\mathbf{B}$
\citep{jackson1999}. The Lorentz equation is a special case of the
partition Lagrangian for an appropriate choice of field topology.
\end{corollary}

\subsection{The four analyser scaling laws}

We specialise to the four canonical mass analyser geometries.

\textbf{Time-of-flight}: $\Depth_{\mathrm{TOF}}(z)=-\kappa z$,
$\Apartition=0$. Solving $\partinert\ddot z=\kappa$ from rest over
length $L$ gives
\begin{equation}
T_{\mathrm{TOF}}=\sqrt{2\partinert L/\kappa}\propto\sqrt{m/z}
\label{eq:tof}
\end{equation}
\citep{wiley1955}.

\textbf{Quadrupole}: $\Depth_{\mathrm{Q}}=\tfrac{\kappa_{0}}{2}(x^{2}-y^{2})[U+V\cos\Omega t]$,
$\Apartition=0$. Substituting $\xi=\Omega t/2$ yields the Mathieu
equation $\ddot u+(a-2q\cos 2\xi)u=0$ with
\begin{equation}
a=\frac{4\kappa_{0}U}{\partinert\Omega^{2}},\quad
q=\frac{2\kappa_{0}V}{\partinert\Omega^{2}}\propto\frac{1}{m/z}
\label{eq:mathieu}
\end{equation}
\citep{paul1990,march2005}.

\textbf{Orbitrap}: $\Depth_{\mathrm{Orb}}=\tfrac{\kappa}{2}(z^{2}-r^{2}/2)
+\tfrac{\kappa R_{m}^{2}}{2}\ln(r/R_{m})$, $\Apartition=0$. The
$z$-equation decouples to $\partinert\ddot z=-\kappa z$, giving
\begin{equation}
\boxed{\;\omega_{z}=\sqrt{\kappa/\partinert}\;}
\quad\propto\sqrt{z/m}
\label{eq:orb}
\end{equation}
\citep{makarov2000}.

\textbf{FT-ICR}: $\Depth=0$, $\Apartition=(B/2)(-y,x,0)$ yields
$\mathbf{B}_{\Depth}=B\hat z$ and
\begin{equation}
\omega_{c}=\frac{qB}{\partinert}=\frac{zeB}{m}\propto z/m
\label{eq:ftcir}
\end{equation}
\citep{marshall1998}.

\begin{theorem}[Analyser Universality]
\label{thm:univ}
Equations \eqref{eq:tof}--\eqref{eq:ftcir} are specialisations of one
Lagrangian \eqref{eq:Lagdef} under choice of partition-field topology.
\end{theorem}

% =========================================================
\section{The Dissolution of Loschmidt's Paradox}
\label{sec:dissolution}
% =========================================================

\subsection{Statement of the paradox}

Loschmidt's paradox is the observation that the microscopic equations
of motion are time-reversal symmetric: if $(\mathbf{q}(t),\mathbf{p}(t))$
is a solution, so is $(\mathbf{q}(-t),-\mathbf{p}(-t))$. Yet macroscopic
processes exhibit time asymmetry: entropy increases. The conventional
formulation reasons: reverse all velocities at time $\tau$; the system
must retrace its trajectory; therefore entropy must \emph{decrease}
back to its initial value at time $2\tau$. Since this is not observed,
something is wrong with the foundation of statistical mechanics.

\subsection{Where the conventional reading goes wrong}

The conventional reading conflates two distinct objects:

\begin{enumerate}[nosep,leftmargin=*]
\item \textbf{The configuration}: the structural arrangement of
correlations, a partition state $\PhiMap(M)$.
\item \textbf{The cumulative configuration count}: $C(t)$, indexed by
$M\in\Zp$, monotone non-decreasing under any traversal direction
(Theorem~\ref{thm:mono}).
\end{enumerate}

When the system at time $2\tau$ occupies a configuration of the same
structural type $(n,\ell,m,s)$ as at time $0$, the conventional reading
identifies these as ``the same state.'' They are not. They are two
distinct occupation events, indexed by distinct integers
$M_{0}<M_{2\tau}\in\Zp$, related by $M_{2\tau}=M_{0}+2N_{\tau}$
(Corollary~\ref{cor:notinitial}). The structural type matches; the
cumulative count does not.

The entropy of the system, in the formulation where entropy is the
configuration count, is $S(t)=\kB C(t)$. By
Theorem~\ref{thm:mono}, $S(t)$ is strictly monotone non-decreasing
in $t$, regardless of velocity-direction choices. Loschmidt's
operation \emph{cannot decrease entropy because the velocity reversal
does not undo the cumulative count}; it traverses configurations of
the same structural types as before, each of which adds to the count.

\subsection{Self-undermining structures in the conventional reading}

\textbf{The metrology argument.} The SI second is defined as
$9{,}192{,}631{,}770$ oscillation periods of caesium-133, each
period containing two velocity reversals of the relevant electronic
motion \citep{si2019}. If reversal cancelled count, this number would
be zero and the unit would be undefined. The unit exists. Therefore
reversal accumulates rather than cancels.

This is a self-undermining structure: Loschmidt's conventional
reading would null the very unit in which the reading is stated. The
reading uses, as its operational primitive, a quantity defined as the
empirical denial of the reading.

\textbf{The pendulum/oscillator argument.} A simple pendulum reverses
its velocity twice per period at the turning points. If reversal
cancelled time-progression, the net period would be zero and pendulums
would not function as clocks. They do function as clocks. The same
argument applies to every physical oscillator, including atomic
clocks, quartz oscillators, and the Orbitrap's axial oscillation.
Each is a velocity reversal performed twice per cycle, accumulating
forward, with the cumulative count grounding metrology.

\subsection{What actually happens}

Under velocity reversal at time $\tau$:

\begin{enumerate}[nosep,leftmargin=*]
\item The system's $(\mathbf{q},\mathbf{p})$ projection retraces. At
time $2\tau$, $(\mathbf{q}_{2\tau},\mathbf{p}_{2\tau})
=(\mathbf{q}_{0},-\mathbf{p}_{0})$. This is the Loschmidt symmetry.
\item The structural type of the configuration at $2\tau$ matches
that at $0$. The same partition cell $(n,\ell,m,s)$ is occupied.
\item The cumulative count $M_{2\tau}$ is strictly greater than $M_{0}$
by twice the number of distinct configurations traversed in $[0,\tau]$.
\item The entropy $S(2\tau)=\kB C(2\tau)>\kB C(0)=S(0)$.
\end{enumerate}

There is no paradox. The Loschmidt symmetry holds for the
$(\mathbf{q},\mathbf{p})$ projection. The entropy increases monotonically
in the cumulative count. These are statements about different
observables; neither contradicts the other.

% =========================================================
\section{Why the Conventional Reading Appears Paradoxical:\\
The Occupation--Specification Asymmetry}
\label{sec:occspec}
% =========================================================

Section~\ref{sec:dissolution} dissolves the paradox structurally. It
does not yet account for why the conventional reading \emph{feels}
incoherent rather than merely wrong. This section gives the diagnosis.
The felt incoherence is the cost an embedded agent must pay to
\emph{specify}---to name, prepare, or recognise---a low-probability
configuration, illegitimately transferred onto the configuration's
\emph{occurrence}, which costs nothing. Three category errors follow
the same transfer pattern; we state them and block each one. The
analysis adds no new physics; it identifies the conceptual move that
makes the dissolution \emph{appear} contested when it is not.

\subsection{Occupation and specification are categorically distinct}

\begin{definition}[Occupation]
\label{def:occupation}
A system \emph{occupies} a microstate $x\in\Omega$ at time $t$ when
the flow has carried it there: $\phi_{t}(x_{0})=x$. Occupation
relates the dynamics to a point. It involves no agent, no choice, and
no preparation.
\end{definition}

\begin{definition}[Specification]
\label{def:specification}
An \emph{embedded agent} of finite descriptive resources
\emph{specifies} a configuration $X\subseteq\Omega$ when it
individuates $X$ from $\Omega\setminus X$ to its finite resolution.
Three modes: (S1) \emph{pointing} (singling $X$ out from its
alternatives), (S2) \emph{preparing} (forcing the system into $X$ by
acting on it), and (S3) \emph{recognising} (certifying that the
system occupies $X$).
\end{definition}

\begin{figure*}[!htbp]
    \centering
    \includegraphics[width=\textwidth]{figures/panel_06_occupation_specification.png}
    \caption{\textbf{Occupation is Free; Specification Pays. The Bedrock Asymmetry.}
    (\textbf{A})~Irrational-rotation trajectory: state $x_{n}$ on the linear $y$-axis ($0$ to $1$) versus step index on the linear $x$-axis ($0$ to $600$), $x_{n+1}=(x_{n}+\alpha)\bmod 1$ with $\alpha=(\sqrt{5}-1)/2$ (the golden mean), $x_{0}=0.314$. The low-measure region $X=[0.10,0.11]$ of measure $p(X)=0.01$ is highlighted (crimson horizontal band). Steelblue markers: every step. Crimson markers: steps that fell inside $X$. The trajectory visits $X$ five times in the first $600$ steps without any agent acting; the visits are unprepared, agent-free occupation events. Empirical content of Proposition~\ref{prop:occfree} (Occupation is Agent-Free).
    (\textbf{B})~Ergodic average: cumulative visit fraction on the linear $y$-axis ($0$ to $0.014$) versus step on a log $x$-axis ($10^{0}$ to $10^{5}$), one trajectory of $10^{5}$ steps. The running average crosses zero, oscillates, and asymptotes to $p(X)=0.01$ (crimson dashed reference) at fractional error $5\times 10^{-16}$ on $200\,000$ samples in the full E7 harness. The visit frequency equals the measure: the only scarce property of $X$ is its measure, not its reachability. Empirical content of Theorem~\ref{thm:rare} clause~(iii) (Birkhoff ergodic average).
    (\textbf{C})~Specification versus occupation cost: cost on the linear $y$-axis (in units of $10^{-20}$~J at $T=300$~K) versus macrostate probability $p(X)$ on a log $x$-axis ($10^{-6}$ to $10^{-1}$). Steelblue: Landauer specification cost $k_{B}T\ln 2\cdot\log_{2}(1/p)$ required to prepare and erase the $\log_{2}(1/p)$ bits identifying $X$. Teal horizontal reference at zero: the cost of \emph{occupying} the same $X$, which is identically zero by Theorem~\ref{thm:rare} clause~(ii). The two costs differ by a divergent positive amount as $p\to 0$. Empirical content of Principle~\ref{prin:asym} (Occupation--Specification Asymmetry) and Theorem~\ref{thm:landauer} (Preparation Pays for Itself).
    (\textbf{D})~Three-dimensional Poincaré return-time densities: histogram density on the vertical axis as a function of return-gap length on the linear $x$-axis ($0$ to $400$ steps) and window width $\mu(X)$ on the linear $y$-axis ($0.005$ to $0.04$), viridis colourmap by width, viewing angle elevation $30^{\circ}$ azimuth $-60^{\circ}$. Four window widths are stacked; in each, the gap distribution is sharply peaked and the mode of the peak shifts toward shorter gaps as $\mu(X)$ grows, $\mathbb{E}[\mathrm{gap}]\propto 1/\mu(X)$, in agreement with the ergodic average of panel~B. Geometric realisation of Theorem~\ref{thm:rare} clause~(i) (recurrence): every window is revisited infinitely often, and the rarity is the measure.}
    \label{fig:panel-occupation-specification}
\end{figure*}

\begin{principle}[Occupation--Specification Asymmetry]
\label{prin:asym}
Let $X$ be a low-probability macrostate. Then:
\begin{enumerate}[nosep,leftmargin=*]
\item \emph{Specification of $X$ is costly}: an embedded agent cannot
single out, prepare, or recognise $X$ without an irreducible
expenditure---boundary residue, imported order, or erased
information---bounded below by a strictly positive constant and, in
the prepared case, no less than the entropy $X$ appears to save.
\item \emph{Occupation of $X$ is free}: a bounded measure-preserving
flow visits $X$ with no agent and no cost; it differs from its visits
to high-probability configurations only in \emph{frequency}.
\end{enumerate}
\end{principle}

The asymmetry is what makes the dissolution sturdy. It is the formal
expression of an everyday observation: a needle in a haystack is
costly to \emph{find} but is in the haystack whether anyone looks,
and the haystack itself contains, at every instant, some specific
and individually improbable arrangement of straws that no one
assembled.

\subsection{Specification is irreducibly costly}

\begin{theorem}[Resolution Floor]
\label{thm:floor}
Let an embedded agent of finite resolution $\delta>0$ individuate a
region $X$ from $\Omega\setminus X$ within a bounded measure space
$\Omega$. Then every realisable separator $\Sigma(X)$ has measure
$\mu(\Sigma(X))\geq\beta>0$, where $\beta$ is the agent's
\emph{resolution floor}. No individuation an embedded agent performs
is free.
\end{theorem}

\begin{proof}
The realisable separator is a finite union of partition cells; by
connectedness it contains at least one cell; by finite resolution
each cell has measure $\geq\mu_{\min}>0$. Set
$\beta:=\mu_{\min}$. Three independent derivations agree on the
constant: geometric (the cell measure above), representational
(any finite description scheme of $\Omega$'s regions has cardinality
$<\mathfrak{c}$ and so cannot exactly name every region; the
unnameable difference has measure $\geq\beta$), and cost-theoretic
(any monotone refinement-cost functional has
$g(t)\to+\infty$ as separator thickness $t\to 0^{+}$, so a
finite-cost process halts at $t_{\star}>0$).
\end{proof}

\begin{theorem}[Preparation Pays for Itself]
\label{thm:landauer}
Let an embedded agent prepare the system into a pre-specified
low-probability macrostate $X$ by acting on it conditional on
knowledge of the system's prior microstate. Let
$\Delta S_{X}=S(W_{\mathrm{typ}})-S(X)>0$ be the entropy deficit of
$X$ relative to a typical macrostate. Then the agent's operation
raises environmental entropy by at least $\Delta S_{X}$:
\[
\Delta S_{\mathrm{env}}\ \geq\ \Delta S_{X}.
\]
The preparation produces no net entropy decrease.
\end{theorem}

\begin{proof}
Conditional preparation requires acquiring $I_{X}\geq\Delta S_{X}/k_{B}$
bits of information distinguishing $X$ from its alternatives. To
reuse a finite memory the agent must erase these bits; by the
Landauer bound \citep{landauer1961,bennett1982}, erasure of $I_{X}$
bits dissipates heat $\geq k_{B}T\ln 2\cdot I_{X}\geq T\Delta S_{X}$,
raising environmental entropy by $\geq\Delta S_{X}$.
\end{proof}

\begin{corollary}[Recognition Cost]
\label{cor:recogcost}
Mode~(S3) recognition deposits boundary residue $\geq\beta$ by
Theorem~\ref{thm:floor}. There is no cost-free recognition.
\end{corollary}

\subsection{Occupation is free; the only scarce property is frequency}

\begin{proposition}[Occupation is Agent-Free]
\label{prop:occfree}
At every instant the system occupies a maximally specific microstate;
no agent specifies, prepares, or recognises it for the occupation to
obtain. The current microstate of any system is, among the microstates
of its macrostate, an individually improbable point that no agent
arranged.
\end{proposition}

\begin{theorem}[Free but Rare Occupation]
\label{thm:rare}
Let $X\subseteq\Omega$ be a macrostate of low probability $p(X)>0$
under the bounded measure-preserving flow. Then:
(i) by Theorem~\ref{thm:poincare}, almost every orbit enters $X$
infinitely often---$X$ is occupied with no agent and no preparation;
(ii) the occupation costs nothing because no agent acts; the flow
merely passes through $X$;
(iii) by the Birkhoff ergodic average, the long-run fraction of time
in $X$ equals $p(X)$, small precisely because $X$ is
low-probability.
\end{theorem}

\begin{proof}
(i) is Poincar\'{e} recurrence
(Theorem~\ref{thm:poincare}) on a finite-measure space; (ii) follows
from Proposition~\ref{prop:occfree} (no actor to incur cost); (iii)
is the indicator-function ergodic average.
\end{proof}

\begin{corollary}[Specificity does not entail Arrangement]
\label{cor:specnotarr}
That a configuration is highly specific (low-probability,
fine-tuned) does not entail that any agent arranged it. The current
microstate is a standing counterexample: maximally specific, arranged
by no one.
\end{corollary}

\subsection{Three category errors in the conventional reading}

The dissolution argument of Section~\ref{sec:dissolution} addresses
two self-undermining structures (metrology, pendulum). We now
identify three further moves the conventional reading uses to
\emph{appear} to keep Loschmidt's paradox alive. Each transfers a
property of specification onto occupation, or silently alters the
measure that constitutes the state. Each is blocked by
Principle~\ref{prin:asym}.

\begin{proposition}[Error A: Renaming a Later State Initial]
\label{prop:errA}
Let a state $y$ be reached by the flow at time $t>0$, so $y$ has a
past. Calling $y$ ``the initial state'' is a labelling (mode~(S1))
and does not endow $y$ with the property---having no past---that
distinguishes an origin. Any theorem whose hypothesis is ``has no
past'' fails to apply to $y$ regardless of $y$'s name.
\end{proposition}

\begin{proof}
``Initial'' admits two senses: (a) a bookkeeping start chosen by the
agent, and (b) the structural origin, the unique state with no past.
The reached $y=\phi_{t}(x_{0})$ has a past by construction; it
satisfies (a) by stipulation but fails (b) by structure. A theorem
proved under (b) reads the predicate ``no past,'' which $y$ lacks.
The rename is inert with respect to that predicate.
\end{proof}

This is exactly Corollary~\ref{cor:notinitial} stated in the agent's
register: $\Phi(M_{2\tau})$ shares the structural type
$(n,\ell,m,s)$ with $\Phi(M_{0})$, but $M_{2\tau}>M_{0}$, so the
later state has a past and cannot inherit the no-past predicate of
the origin.

\begin{proposition}[Error B: Inferring Arrangement from Specificity]
\label{prop:errB}
From ``$X$ is highly specific (low-probability)'' one cannot infer
``an agent arranged $X$.'' The inference confuses a property of the
outcome (its measure) with a claim about its causal origin (an
agent's action).
\end{proposition}

\begin{proof}
Corollary~\ref{cor:specnotarr} exhibits a standing counterexample:
the system's current microstate is maximally specific and arranged
by no one. Specificity is measure; arrangement is agency; the first
does not imply the second.
\end{proof}

\begin{proposition}[Error C: Demanding a Low-Entropy State be Most Probable]
\label{prop:errC}
There is no operation that makes a low-entropy macrostate the most
probable next state while it remains low-entropy. ``Most probable''
and ``high entropy'' are the same predicate up to an additive
constant: for a macrostate $W\subseteq\Omega$, $S(W)=k_{B}\log\mu(W)$
and $\log p(W)=\log\mu(W)-\log\mu(\Omega)$ differ by the constant
$-\log\mu(\Omega)$. Demanding a low-entropy state be most probable
demands $\mu(X)=\min$ and $\mu(X)=\max$ simultaneously---a
contradiction.
\end{proposition}

\begin{remark}
\label{rem:nesting}
Conditioning---``state $a$ makes $b$ probable, $b$ makes $c$
probable''---cannot raise a small target's measure. The share of any
predecessor's onward orbits entering a small region is bounded by the
region's measure; nested chains drive occupation toward larger
macrostates (higher entropy), never toward smaller ones. The nesting
is the engine of the arrow, not a route around it.
\end{remark}

\begin{theorem}[The Three Errors Share One Form]
\label{thm:oneform}
Errors A, B, and C are instances of a single illegitimate transfer:
a property established for the agent's relation to a state
(specification: naming, locating, constructing) is carried onto the
system's relation to the state (occupation), or the constitutive
measure of the state is silently altered. Under
Principle~\ref{prin:asym} and the entropy--probability identity each
transfer is blocked.
\end{theorem}

\subsection{Diagnosis of the felt incoherence}

\begin{theorem}[Spontaneous Reversal is Atypical, not Incoherent]
\label{thm:atypical}
The spontaneous reversed orbit is a legal, agent-free occupation
event (Theorem~\ref{thm:rare}): it exists by recurrence, costs
nothing because no agent acts, and is merely rare because its
macrostate has small measure. Assuming its existence yields no
contradiction---only an improbability of order
$p=\mu(X)/\mu(\Omega)$.
\end{theorem}

\begin{proof}
By Theorem~\ref{thm:rare}, the low-measure reflected macrostate is
occupied infinitely often by almost every orbit, with no agent and
no cost. A contradiction would follow only if existence entailed an
impossible state of affairs; it does not---the orbit is an ordinary
forward trajectory of $\phi_{t}$, distinct from the origin
(Corollary~\ref{cor:notinitial}), depositing positive individuation
residue at each step it is \emph{analysed} (Theorem~\ref{thm:floor})
yet itself costing nothing to \emph{occur}
(Proposition~\ref{prop:occfree}). The felt incoherence is the cost
of \emph{specifying} the orbit (Theorem~\ref{thm:landauer}),
transferred by Errors~B/C onto its \emph{occurrence}. Removing the
transfer, what remains is rarity, not incoherence.
\end{proof}

\begin{remark}[The bedrock fact]
\label{rem:bedrock}
A low-probability configuration is real (it is occupied), uncompensated
(no agent pays), unarranged (no specifier), and rare. Its rarity is
its identity: it cannot be made common without ceasing to be
low-probability (the entropy--probability identity), made impossible
without contradicting recurrence (Theorem~\ref{thm:poincare}), or
made arranged without an agent and a cost
(Theorem~\ref{thm:landauer}). Rarity is the residue that survives
every reframing. The dissolution of
Section~\ref{sec:dissolution} is the same fact stated at the level
of the cumulative count: the rarity is the measure of the macrostate;
the monotone count records its occupations one by one.
\end{remark}

\subsection{The four Loschmidt claims, classified}

\begin{table}[H]
\centering
\small
\begin{tabular}{p{4.8cm}cp{6.8cm}}
\toprule
Claim & Status & Reason \\
\midrule
Reversed endpoint $\equiv$ initial state &
\textbf{false} &
reached forward, has a past;\newline rename is Error A
(Cor.~\ref{cor:notinitial}, Prop.~\ref{prop:errA}) \\
Prepared reversal decreases entropy &
\textbf{compensated} &
Landauer cost $\geq\Delta S$
(Thm.~\ref{thm:landauer}) \\
Spontaneous reversal is incoherent &
\textbf{false} &
occupation is free; only rare\newline
(Thm.~\ref{thm:atypical}) \\
Macrostate entropy falls spontaneously &
\textbf{true, atypical} &
uncompensated but $p(X)\ll 1$\newline
(Thm.~\ref{thm:rare}(iii)) \\
\bottomrule
\end{tabular}
\caption{The four Loschmidt claims under the
occupation--specification distinction. The reversal is coherent and
real; its only scarce property is frequency.}
\label{tab:lossummary}
\end{table}

\subsection{Implications for the headline experiment}

The Cs\textsuperscript{+} Orbitrap experiment of
Section~\ref{sec:experiment} tests Theorem~\ref{thm:rare} and
Theorem~\ref{thm:landauer} simultaneously. The trap performs
$\sim 10^{6}$ velocity reversals per second of analytical
observation: each is an unprepared, agent-free traversal of a
configuration whose structural type resembles an earlier one. The
Hilbert-transform receiver is structurally distinct from the orbit
(it lives in the readout chain, not in the trap), so the
recognition gap (Section~\ref{sec:backward}) is not closed by the
measurement; the cumulative count is measured \emph{about} the
orbit, not from inside it. The pass criteria of
Protocol~\ref{prot:C1} are therefore tests of agent-free
monotonicity: the orbit occupies its trajectory at no cost; the
receiver pays a finite---and bounded---specification cost to record
that occupation; the two are compatible because they obey opposite
economies (Principle~\ref{prin:asym}).

\begin{figure*}[!htbp]
    \centering
    \includegraphics[width=\textwidth]{figures/panel_03_hilbert.png}
    \caption{\textbf{Hilbert-Transform Phase Recovery and the Measured Partition Count.}
    (\textbf{A})~Noisy transient $s(t)$ on the linear $y$-axis ($-1$ to $+1$) versus time on the linear $x$-axis ($0$ to $200$~ms), synthesised as $\cos(\omega t)$ at $\omega/(2\pi)=50$~Hz with additive Gaussian noise of standard deviation $0.02$ and sampled at $f_{s}=5000$~Hz ($N=1000$ samples). Ten complete cycles are visible; the noise is sub-percent. This is the input to the analytic-signal construction used downstream.
    (\textbf{B})~Hilbert envelope: instantaneous amplitude $|s(t)+i\mathcal{H}[s(t)]|$ on the linear $y$-axis ($-1$ to $+1$) versus time on the linear $x$-axis ($0$ to $200$~ms). Faint cyan trace: the noisy signal. Crimson curves: $\pm|s+i\mathcal{H}s|$, the upper and lower envelope. The envelope is constant at unity to within noise; the analytic signal lives on the unit circle. This is the precondition for unambiguous phase recovery.
    (\textbf{C})~Phase-recovery residual $M^{\mathrm{meas}}(t)-M^{\mathrm{pred}}(t)$ on the linear $y$-axis (in units of $10^{-3}$) versus time on the linear $x$-axis ($0$ to $200$~ms), with $M^{\mathrm{meas}}(t)=\theta_{\mathrm{unwrap}}(t)/(2\pi)$ and $M^{\mathrm{pred}}(t)=\omega t/(2\pi)$ (neutral horizontal reference at zero). The residual is bounded by $|M^{\mathrm{meas}}-M^{\mathrm{pred}}|<0.01$ over the whole transient; full-length harness run E4 gives a relative error of $1.6\times 10^{-11}$ at $T=1$~s over $10^{6}$ cycles. Empirical content of Corollary~\ref{cor:phase} (Exact Phase) under finite-sample, noisy conditions.
    (\textbf{D})~Three-dimensional analytic-signal trace: $\mathrm{Re}(s+i\mathcal{H}s)$ on the linear $x$-axis, $\mathrm{Im}(s+i\mathcal{H}s)$ on the linear $y$-axis, time on the vertical axis ($0$ to $200$~ms), viewing angle elevation $30^{\circ}$ azimuth $-60^{\circ}$. The trace is a perfect helix at constant radius, winding monotonically up the $t$-axis. This is the geometric realisation of the time-count identity $\theta(M)=2\pi M$ (Corollary~\ref{cor:phase}): each loop of the helix is one increment of $M$, and the helix never reverses.}
    \label{fig:panel-hilbert}
\end{figure*}

% =========================================================
\section{Backward Trajectory Completion: The Penultimate State}
\label{sec:backward}
% =========================================================

\subsection{The S-entropy embedding}

\begin{definition}[S-Entropy Coordinates]
\label{def:Sspace}
The S-entropy coordinate space is $\Sspace=[0,1]^{3}$ with coordinates
$(S_{k},S_{t},S_{e})$ representing structural, temporal, and
evolutionary projections of the configuration count.
\end{definition}

\begin{definition}[Ternary Refinement Hierarchy]
\label{def:hier}
A ternary refinement hierarchy of depth $k$ in $\Sspace$ is a sequence
of partitions $\{P_{0},\ldots,P_{k}\}$ with $|P_{0}|=1$,
$|P_{j+1}|=3|P_{j}|$, each block of $P_{j}$ refined into three
children. Total leaves: $N=3^{k}$.
\end{definition}

\subsection{Backward navigation in $\log_{3}N$ steps}

\begin{theorem}[Backward Navigation]
\label{thm:logn}
Given the endpoint cell $\tau_{k}\in\{0,1,2\}^{k}$, backward navigation
to the root terminates in exactly $\log_{3}N=k$ steps.
\end{theorem}

\begin{proof}
Each step computes the parent of the current cell by truncating the
last trit of its base-3 expansion. This is $O(1)$ per step. Starting
at depth $k$ and ending at depth $0$, exactly $k$ steps are required.
\end{proof}

\begin{algorithm}[!htbp]
\caption{Backward Trajectory Completion}
\label{alg:back}
\begin{algorithmic}[1]
\Function{Complete}{addr $\tau\in\{0,1,2\}^{k}$}
\State trajectory $\mathcal{T}\gets[\tau]$
\State $j\gets k$
\While{$j>1$}
\State $\tau\gets\tau[1{:}j-1]$
\Comment{drop last trit (parent)}
\State $\mathcal{T}.\mathrm{prepend}(\tau)$
\State $j\gets j-1$
\EndWhile
\State \Return $\mathcal{T}$
\Comment{terminates at the penultimate state}
\EndFunction
\end{algorithmic}
\end{algorithm}

\subsection{Termination at the penultimate state}

Algorithm~\ref{alg:back} terminates the \texttt{while} loop at $j=1$,
producing the penultimate state---one isomorphism step short of the
root. This is not an off-by-one artefact but a structural requirement.

\begin{theorem}[Recognition Gap]
\label{thm:recog}
Recognition of a configuration $\tau\in\{0,1,2\}^{k}$ by a receiver
requires the receiver's state to be distinct from $\tau$. Closing the
final step---making the receiver state identical to $\tau$---dissolves
the distinction and with it the recognition. The minimum distinct
vantage is bounded below by the resolution floor: any separator
between receiver and $\tau$ has measure $\geq\beta>0$
(Theorem~\ref{thm:floor}), so the gap is not an artefact of choice
but a structural lower bound.
\end{theorem}

\begin{proof}
A receiver that occupies $\tau$ does not stand in relation to $\tau$
as a recognising party to a recognised object; the two are
indistinguishable, and recognition (mode~(S3),
Definition~\ref{def:specification}) requires individuation against
$\Omega\setminus\{\tau\}$. By Theorem~\ref{thm:floor}, the realisable
separator has measure $\geq\beta$; the receiver therefore stands at
distance $\geq\beta$ from $\tau$ in the agent's partition. The
minimal distinct vantage is the penultimate state, one isomorphism
step from $\tau$. Closing that gap would reduce separator measure to
zero, contradicting Theorem~\ref{thm:floor}.
\end{proof}

\begin{remark}[Self-referential verification]
\label{rem:selfref}
The deeper reason the gap cannot close is self-referential. To
\emph{verify} that the receiver occupies $\tau$ is to decide its
membership in the collection that constitutes $\tau$ (the negation
of everything $\tau$ is not). Self-referential set-membership has
no consistent value (Russell, Cantor, Tarski). The gap is therefore
not merely a finite-resolution artefact: it is the geometric
shadow of a logical obstruction. Theorem~\ref{thm:floor} bounds the
shadow's measure; the obstruction itself is not exhibitable as a set.
\end{remark}

\begin{figure*}[!htbp]
    \centering
    \includegraphics[width=\textwidth]{figures/panel_04_resolution_floor.png}
    \caption{\textbf{The Resolution Floor: Three Independent Derivations Agreeing on One Constant.}
    (\textbf{A})~Floor $\beta$ on a log $y$-axis ($10^{-6}$ to $10^{0}$) versus partition scale $h$ on a log $x$-axis ($10^{-3}$ to $10^{-1}$), $40$ logarithmically spaced $h$ values. Steelblue: geometric floor $\beta_{\mathrm{geom}}=h^{2}$ from the minimum cell measure in $[0,1]^{2}$. Crimson: representational residue $\beta_{\mathrm{rep}}=h^{2}\max(1,2\pi r/h)$ for an irrational-radius disc $r=1/\sqrt{2}$. Both are strictly positive at every tested $h$; both vanish only in the excluded limit $h\to 0$. Empirical content of Theorem~\ref{thm:floor} (Resolution Floor): the geometric and representational floors bound the same constant from below away from zero.
    (\textbf{B})~Cost functional $g(t)=-\log t$ on the linear $y$-axis ($0$ to $5$) versus separator thickness $t$ on the linear $x$-axis ($0.005$ to $1$). The curve diverges as $t\to 0^{+}$; a fixed cost budget $C=5$ (neutral dashed horizontal reference) intersects the curve at $t_{\star}=e^{-5}\approx 6.74\times 10^{-3}$ (teal dashed vertical reference). No finite-cost process attains a sharp separator. Empirical content of the cost derivation of Theorem~\ref{thm:floor}: the cost-bounded thickness $t_{\star}$ is strictly positive for any finite budget.
    (\textbf{C})~Agreement ratio $\beta_{\mathrm{rep}}/\beta_{\mathrm{geom}}$ on the linear $y$-axis ($0$ to $4500$) versus partition scale $h$ on a log $x$-axis ($10^{-3}$ to $10^{-1}$). The ratio is bounded above and grows monotonically as $h\to 0$, with the representational floor exceeding the geometric floor by the irrational-disc perimeter factor $2\pi r/h$. The two derivations agree in sign and order; neither can be reduced to zero while the other remains positive. Empirical content of Theorem~\ref{thm:floor} (agreement clause).
    (\textbf{D})~Three-dimensional floor surface: $\log_{10}\beta$ on the vertical axis as a function of $\log_{10} h$ on the linear $x$-axis ($-3$ to $-1$) and cost budget $C$ on the linear $y$-axis ($1$ to $12$), viridis colourmap, viewing angle elevation $30^{\circ}$ azimuth $-60^{\circ}$. The surface is $\log_{10}\max(h^{2},e^{-C})$: the floor is governed by the larger of the geometric and cost contributions at each $(h,C)$. The surface lies strictly above $-\infty$ everywhere; the floor is bounded below by a positive constant on every section. Geometric realisation of Theorem~\ref{thm:floor} (the constant is one).}
    \label{fig:panel-floor}
\end{figure*}

\begin{corollary}[Loschmidt's Operation Cannot Close the Gap]
\label{cor:losch_gap}
Loschmidt's velocity reversal asks whether the cumulative count can
return to its initial value. By Theorem~\ref{thm:recog}, closing the
gap would dissolve the distinction between the reversed state and the
initial state; by Corollary~\ref{cor:notinitial}, this does not happen
under velocity reversal because the cumulative count strictly
increases. The penultimate-state structure is the operational analogue
of the impossibility of returning to the initial cumulative count.
\end{corollary}

\subsection{Virtual sub-states and the speedup}

\begin{definition}[Sub-coordinate Decomposition]
\label{def:sub}
A sub-coordinate decomposition of $s\in[0,1]$ is a triple
$(s_{1},s_{2},s_{3})\in\R^{3}$ with $(s_{1}+s_{2}+s_{3})/3=s$. It is
\emph{virtual} if any $s_{i}\notin[0,1]$.
\end{definition}

\begin{theorem}[Virtual Decompositions Dominate]
\label{thm:virt}
For any $s\in(0,1)$ and bound $M$, the fraction of admissible
decompositions $(s_{1},s_{2},s_{3})$ with $\max_{i}|s_{i}|\leq M$
that are virtual tends to $1$ as $M\to\infty$.
\end{theorem}

\begin{proof}
The constraint plane $\{(s_{1},s_{2},s_{3}):s_{1}+s_{2}+s_{3}=3s\}$
is two-dimensional in $\R^{3}$. Its intersection with $[0,1]^{3}$ has
bounded area; its intersection with $[-M,M]^{3}$ has area growing as
$M^{2}$. The complement (virtual) fraction approaches $1$.
\end{proof}

\begin{theorem}[Collapse Without Virtual Decompositions]
\label{thm:collapse}
Backward completion restricted to physical decompositions
($s_{i}\in[0,1]$ for all $i$) requires $\Theta(N)$ steps.
\end{theorem}

\begin{proof}
A physical-only step requires the metric to be evaluated at resolution
finer than the cell width at the parent level; this is achievable only
via decompositions whose components escape $[0,1]^{3}$
(Theorem~\ref{thm:virt}). Without virtual decompositions the parent
must be found by enumerating children of candidate parents, summing
to $\sum_{j=0}^{k}3^{j}=\Theta(N)$. The same conclusion follows from
the cost derivation of Theorem~\ref{thm:floor}: a monotone
refinement-cost functional $g$ satisfies $g(t)\to+\infty$ as
separator thickness $t\to 0^{+}$; restricting to physical
decompositions caps $t$ at the cell width at depth $k$, forcing
either divergent cost per step or enumeration of all $3^{k}$ cells.
\end{proof}

% =========================================================
\section{Phase Coherence and Subharmonic Encoding}
\label{sec:phase}
% =========================================================

\subsection{Phase relations}

\begin{theorem}[Precursor--Fragment Phase Coherence]
\label{thm:coh}
Let a precursor ion at partition count $M_{\mathrm{prec}}$ be on a
continuous oscillatory trajectory. Let $f\in\calF(M_{\mathrm{prec}})$
be a member of its forward sequence at count
$M_{f}>M_{\mathrm{prec}}$. The phase relationship satisfies
\begin{equation}
\boxed{\;\Delta\theta_{f}=2\pi(M_{f}-M_{\mathrm{prec}}).\;}
\label{eq:dtheta}
\end{equation}
\end{theorem}

\begin{proof}
By Corollary~\ref{cor:phase}, $\theta(M)=2\pi M$. Subtracting the
phases at $M_{f}$ and $M_{\mathrm{prec}}$ gives \eqref{eq:dtheta}.
\end{proof}

\subsection{Fragment frequencies and subharmonics}

\begin{theorem}[Fragment Frequency Law]
\label{thm:ffreq}
In an Orbitrap with parameter $\kappa$, a fragment ion $f$ of mass
$m_{f}$ derived from a precursor of mass $m_{\mathrm{prec}}$ (same
charge, same $\kappa$) oscillates at axial frequency
\begin{equation}
\omega_{f}=\omega_{\mathrm{prec}}\sqrt{m_{\mathrm{prec}}/m_{f}}.
\label{eq:ffreq}
\end{equation}
\end{theorem}

\begin{proof}
From \eqref{eq:orb}, $\omega^{2}\propto 1/m$ at fixed charge and
$\kappa$. Take the ratio.
\end{proof}

\begin{theorem}[Subharmonic Encoding]
\label{thm:sub}
For a precursor ion with Orbitrap transient
$s_{\mathrm{prec}}(t)=A_{\mathrm{prec}}\cos(\omega_{\mathrm{prec}}t+\phi_{\mathrm{prec}})$,
each fragment $f$ contributes an additive component
\begin{equation}
s_{f}(t)=A_{f}\cos(\omega_{f}t+\phi_{f}),
\label{eq:sub}
\end{equation}
with $\omega_{f}$ from \eqref{eq:ffreq} and
$\phi_{f}=\phi_{\mathrm{prec}}-2\pi(M_{f}-M_{\mathrm{prec}})$
(Theorem~\ref{thm:coh}). The Fourier transform of the precursor
transient contains spectral peaks at $\omega_{f}$ for every fragment
in the forward sequence reachable during the transient.
\end{theorem}

\begin{proof}
The fragment acquires a distinct oscillator identity at the moment of
fragmentation; from that moment, its image current is additive with
the precursor's by linearity of charge-induced detection. Its initial
phase is set by the precursor's phase at the fragmentation moment
minus the phase advance during $[0,\tau_{f}]$, which by
Theorem~\ref{thm:coh} is $2\pi\Delta M_{f}$.
\end{proof}

\begin{remark}[Distinguishability from Harmonics]
\label{rem:harm}
Fragment subharmonics appear at irrational multiples
$\sqrt{m_{\mathrm{prec}}/m_{f}}$ of the precursor frequency, distinct
from integer-multiple harmonics arising from trap nonlinearity. The
two are distinguishable in any FFT of adequate bin resolution.
\end{remark}

% =========================================================
\section{The Headline Experiment: Caesium in an Orbitrap}
\label{sec:experiment}
% =========================================================

\subsection{Rationale}

The Cs\textsuperscript{+} ion in an Orbitrap mass spectrometer is the
optimal operational instantiation of the framework's central claim,
for four reasons.

\textbf{Reason 1: The analyte is the metrology atom.} Caesium-133 is
the isotope whose hyperfine count defines the SI second
\citep{si2019}. The experiment uses, as its test object, the very
atom whose monotonic accumulation of velocity-reversal counts grounds
international metrology.

\textbf{Reason 2: The trap performs the reversal natively.} The
Orbitrap's axial well drives the ion through approximately
$\omega_{z}T/(2\pi)\sim 10^{6}$ axial oscillations per second of
analytical observation, each containing two velocity reversals at the
turning points. The trap is, by construction, a controlled-reversal
apparatus. No external reversal mechanism is required.

\textbf{Reason 3: The Lagrangian gives an exact prediction.} From
\eqref{eq:orb} and \eqref{eq:tc}, the partition count is
\begin{equation}
M^{\mathrm{pred}}(T)=\frac{\omega_{z}T}{2\pi}=\frac{T}{2\pi}\sqrt{\frac{e\kappa}{m_{\mathrm{Cs}}}},
\label{eq:Mpred}
\end{equation}
with $m_{\mathrm{Cs}}=132.905452\,\mathrm{u}$ \citep{atomicmasses2020}
and $\kappa$ a trap geometry constant. There are no adjustable
parameters.

\textbf{Reason 4: The transient gives an exact measurement.} The
Orbitrap image current $s(t)$ has instantaneous phase
$\theta(t)=\arg(s(t)+i\mathcal{H}[s(t)])$ recovered by Hilbert
transform, where $\mathcal{H}$ is the Hilbert operator \citep{gabor1946}.
Cumulative unwrapped phase divided by $2\pi$ yields a measured
partition count
\begin{equation}
M^{\mathrm{meas}}(t)=\theta_{\mathrm{unwrap}}(t)/(2\pi)
\label{eq:Mmeas}
\end{equation}
at every sample.

\subsection{Protocol}

\begin{protocol}[Caesium Monotonicity Experiment]
\label{prot:C1}
\textbf{Input.} Orbitrap raw transient of Cs\textsuperscript{+}
($m/z\approx 132.9054$, singly charged) recorded over duration
$T\in[0.5,2.0]$~s with trap parameter $\kappa$ from instrument
metadata. Sampling rate $f_{s}\geq 4\omega_{z}/(2\pi)$ to ensure
adequate Nyquist coverage.

\textbf{Predicted statistic.}
$M^{\mathrm{pred}}(T)=(T/2\pi)\sqrt{e\kappa/m_{\mathrm{Cs}}}$ from
\eqref{eq:Mpred}.

\textbf{Measured statistic.} $M^{\mathrm{meas}}(t)$ from
\eqref{eq:Mmeas}, sampled at the ADC rate.

\textbf{Pass criterion 1 (Lagrangian agreement).}
\[
\frac{|M^{\mathrm{meas}}(T)-M^{\mathrm{pred}}(T)|}{M^{\mathrm{pred}}(T)}<10^{-6}.
\]

\textbf{Pass criterion 2 (Monotonicity).}
$M^{\mathrm{meas}}(t_{i+1})>M^{\mathrm{meas}}(t_{i})$ for every
adjacent sample pair $(t_{i},t_{i+1})$ over the entire transient.
\end{protocol}

\subsection{What this falsifies}

\textbf{Conventional Loschmidt prediction.} Velocity reversal at each
turning point should cause $M(t)$ to decrement (retracing of the
counting process). For an Orbitrap transient of duration
$T=1.0$~s, this should occur approximately
$2\omega_{z}/(2\pi)\sim 10^{6}$ times. Pass criterion 2 in
Protocol~\ref{prot:C1} is the explicit denial of this prediction.

\textbf{Framework prediction.} $M(t)$ is strictly monotone
non-decreasing across all turning points; agreement with
$M^{\mathrm{pred}}(T)$ to fractional precision $10^{-6}$.

If $M^{\mathrm{meas}}(t)$ exhibits any decrement at any turning point,
the framework is operationally falsified. If $M^{\mathrm{meas}}(t)$ is
monotone over $\sim 10^{6}$ velocity reversals and matches the
Lagrangian prediction to $10^{-6}$, the Loschmidt operation has been
performed $\sim 10^{6}$ times on the SI metrology atom, in a single
analytical observation, and observed to leave the cumulative count
monotone every time.

\subsection{Implementation notes}

The Hilbert transform of the discrete-time transient is computed via
FFT \citep{cooley1965}: take the FFT, zero the negative-frequency
components, double the positive-frequency components, inverse FFT,
and take the imaginary part. Phase unwrapping uses the standard
algorithm: detect discontinuities exceeding $\pi$ in successive phase
samples and add multiples of $2\pi$ as needed. The trap parameter
$\kappa$ is read from instrument configuration; alternatively, it can
be calibrated against a known reference ion in the same transient
without affecting the monotonicity test.

% =========================================================
\section{Supporting Experiments}
\label{sec:supporting}
% =========================================================

The headline experiment of Section~\ref{sec:experiment} is the
primary operational test. The following four supporting experiments
provide auxiliary validation of the framework's structural results on
existing data, with no new wet-lab work required.

\subsection{Bijection on standard \texttt{mzML} ions}

\begin{protocol}[Bijection Check]
\label{prot:bij}
\textbf{Input.} \texttt{mzML} file with $N_{\mathrm{ion}}$ detected
ions, each with $m/z$ and intensity.

\textbf{Procedure.} For each ion, compute
$n_{i}=\lfloor\sqrt{(m/z)_{i}/m_{\mathrm{ref}}}\rfloor+1$ for
$m_{\mathrm{ref}}=1$~Da, encode the ion as
$M_{i}=N_{\mathrm{state}}(n_{i}-1)+\delta_{i}$ with $\delta_{i}$
constructed from the ion's secondary observables, apply $\PhiMap$ to
recover $(n_{i},\ell_{i},m_{i},s_{i})$, and verify capacity
$\delta_{i}\leq 2n_{i}^{2}$ and selection $|m_{i}|\leq\ell_{i}<n_{i}$.

\textbf{Pass criterion.} All $N_{\mathrm{ion}}$ ions satisfy capacity
and selection; no two distinct ions share an integer encoding.

\textbf{What this validates.} Theorems~\ref{thm:cap}, \ref{thm:bij},
and Corollary~\ref{cor:cum}.
\end{protocol}

\subsection{Four analyser scaling laws}

\begin{protocol}[Analyser Scaling]
\label{prot:scale}
\textbf{Input.} Reference dataset of $N_{c}\geq 30$ compounds with
measured flight times (TOF), Orbitrap frequencies (Orb), Mathieu
$(a,q)$ (Quad), and cyclotron frequencies (FT-ICR).

\textbf{Procedure.} For each analyser, fit the predicted scaling
\eqref{eq:tof}--\eqref{eq:ftcir} and compute the maximum relative
residual across compounds.

\textbf{Pass criterion.} All four maximum residuals
$<10^{-4}$ on a dataset spanning at least one decade in $m/z$.

\textbf{What this validates.} Theorem~\ref{thm:univ}: the four
analyser equations are specialisations of one Lagrangian.
\end{protocol}

\subsection{Subharmonic encoding}

\begin{protocol}[Single-Transient Fragment Encoding]
\label{prot:sub}
\textbf{Input.} Raw Orbitrap transient $s(t)$ of one precursor at
$\omega_{\mathrm{prec}}$, paired with the corresponding centroided
MS\textsuperscript{2} fragment list
$\{(m_{f_{i}},I_{f_{i}})\}_{i=1}^{N_{f}}$.

\textbf{Procedure.} Compute predicted fragment frequencies
$\omega_{f_{i}}^{\mathrm{pred}}=\omega_{\mathrm{prec}}\sqrt{m_{\mathrm{prec}}/m_{f_{i}}}$
and predicted phases
$\phi_{f_{i}}^{\mathrm{pred}}=\phi_{\mathrm{prec}}-2\pi\Delta M_{f_{i}}$.
Compute $\hat s(\omega)$ by FFT with bin width
$\Delta\omega<\min_{i}|\omega_{f_{i}}-\omega_{\mathrm{prec}}|/10$.
For each predicted $\omega_{f_{i}}^{\mathrm{pred}}$, search $\hat s$
within $\pm 3\Delta\omega$ for a peak above SNR $5$. Record measured
frequency and phase residuals.

\textbf{Pass criterion.} For $\geq 80\%$ of MS\textsuperscript{2}
fragments above SNR floor: frequency residual $<10^{-5}$, phase
residual $<0.1$~rad modulo $\pi$.

\textbf{What this validates.} Theorems~\ref{thm:coh}, \ref{thm:ffreq},
\ref{thm:sub}.
\end{protocol}

\subsection{Backward completion complexity}

\begin{protocol}[Logarithmic Completion]
\label{prot:back}
\textbf{Input.} $N_{c}\geq 30$ ions with annotated S-entropy
coordinates $(S_{k},S_{t},S_{e})$ at depth $k=18$.

\textbf{Procedure (virtual).} Apply Algorithm~\ref{alg:back}; record
step count $K_{i}^{\mathrm{virt}}$.

\textbf{Procedure (physical-only).} Apply Algorithm~\ref{alg:back}
restricted to decompositions with $s_{j}\in[0,1]$; record step count
$K_{i}^{\mathrm{phys}}$.

\textbf{Pass criterion.} $K_{i}^{\mathrm{virt}}=k-1=17$ exactly (the
algorithm terminates at the penultimate state) for all $i$;
$K_{i}^{\mathrm{phys}}$ scales as $3^{k}$ within a constant factor.

\begin{figure*}[!htbp]
    \centering
    \includegraphics[width=\textwidth]{figures/panel_05_backward.png}
    \caption{\textbf{Backward Completion Complexity and the Penultimate-State Halt.}
    (\textbf{A})~Steps to root on a log $y$-axis ($10^{0}$ to $10^{8}$) versus tree depth $k$ on the linear $x$-axis ($1$ to $17$). Steelblue circles: virtual-decomposition steps $K^{\mathrm{virt}}=k-1$, growing linearly in $k$. Crimson squares: physical-decomposition step count $K^{\mathrm{phys}}=(3^{k+1}-1)/2$, growing geometrically. At $k=15$ the physical algorithm requires $2.15\times 10^{7}$ steps; the virtual algorithm requires $14$. Empirical content of Theorem~\ref{thm:logn} (Backward Navigation) and Theorem~\ref{thm:collapse} (Collapse Without Virtual Decompositions).
    (\textbf{B})~Speedup ratio $K^{\mathrm{phys}}/K^{\mathrm{virt}}$ on a log $y$-axis ($10^{0}$ to $10^{8}$) versus depth $k$ on the linear $x$-axis ($1$ to $17$). The ratio grows as $\sim 3^{k+1}/(2k)$, exceeding $10^{6}$ by $k=14$. The virtual-substate decomposition is not a constant-factor improvement: it is an asymptotic separation between $\log_{3}N$ and $\Theta(N)$.
    (\textbf{C})~Penultimate-state certificate: observed virtual step count on the linear $y$-axis versus depth $k$ on the linear $x$-axis ($3$ to $15$), error bars showing standard deviation across $40$ random endpoints per depth (crimson dashed reference at $k-1$). Mean and max and min all coincide with $k-1$ at every depth; standard deviation is zero. Algorithm~\ref{alg:back} terminates at exactly the penultimate state for every sampled endpoint, never closing the final step. Empirical content of Theorem~\ref{thm:recog} (Recognition Gap).
    (\textbf{D})~Three-dimensional ternary-tree cloud at $k=5$ ($N=243$ leaves) embedded in S-entropy space $(S_{k},S_{t},S_{e})\in[0,1]^{3}$ by interpreting successive trits as base-$3$ refinements on alternating axes, viewing angle elevation $30^{\circ}$ azimuth $-60^{\circ}$. Marker colour encodes the leaf norm $\|S\|$ on the viridis scale. The cloud is the discrete trajectory cone of all $3^{k}$ endpoints; backward navigation from any one of them reaches the penultimate level in $k-1$ steps and halts. Geometric realisation of Definition~\ref{def:hier} (Ternary Refinement Hierarchy) and the algorithm's termination structure.}
    \label{fig:panel-backward}
\end{figure*}

\textbf{What this validates.} Theorems~\ref{thm:logn}, \ref{thm:collapse}.
\end{protocol}

% =========================================================
\section{JSON Record Schemas and Suite-Level Acceptance}
\label{sec:json}
% =========================================================

\subsection{Common record schema}

Each experiment persists a JSON record:
\begin{verbatim}
{
  "experiment": "C1|Bij|Scale|Sub|Back",
  "theorem_ids": [<int>, ...],
  "input_dataset": <str>,
  "n_samples": <int>,
  "predicted": { ... },
  "measured": { ... },
  "residuals": {
    "max": <float>,
    "rms": <float>
  },
  "monotone": <bool>,
  "pass": <bool>,
  "elapsed_seconds": <float>,
  "seed": <int>
}
\end{verbatim}

The aggregate file \texttt{master\_results.json} collects per-experiment
pass flags, total wall-clock time, and the framework version.

\subsection{Suite-level acceptance}

The full validation suite is accepted if and only if:
\begin{itemize}[nosep,leftmargin=*]
\item Protocol~\ref{prot:C1} (headline) passes both criteria.
\item Each supporting protocol passes its criterion.
\end{itemize}

A \texttt{FAIL} on the headline protocol falsifies the framework
operationally: the conventional Loschmidt prediction
($M$ retraces at turning points) would have been observed. A
\texttt{FAIL} on a supporting protocol localises the failure to the
specific theorem under test.

% =========================================================
\section{Discussion}
\label{sec:discussion}
% =========================================================

\subsection{What this paper proves and what it does not}

The derivations of Sections~\ref{sec:axiom}--\ref{sec:phase} establish
that, granted Axiom~\ref{ax:bps} alone, the configuration count
$C(t)$ is well defined and monotone non-decreasing under any
traversal direction, including velocity reversal. They establish that
Loschmidt's conventional reading conflates two distinct objects:
the structural type of the configuration (which can recur), and the
cumulative count of occupation events (which cannot).

The paper does not establish, by derivation alone, that the framework
applies physically to caesium ions in Orbitraps. That is the role of
the headline experiment of Section~\ref{sec:experiment}. If the
experiment fails---if $M^{\mathrm{meas}}$ exhibits decrements at
turning points, or fails to match $M^{\mathrm{pred}}$ to $10^{-6}$---the
framework is falsified at the operational level.

\subsection{Limitations}

\begin{enumerate}[nosep,leftmargin=*]
\item \textbf{Integrable systems.} The Lagrangian derivation assumes
integrability. KAM-type robustness \citep{kolmogorov1954,arnold1963,moser1962}
extends results to near-integrable cases; chaotic regimes require
additional treatment.
\item \textbf{Detector noise.} The Hilbert-transform phase recovery
assumes signal-to-noise sufficient to track phase between samples.
Practical Orbitrap transients are well above this floor for the
duration $T$ specified in Protocol~\ref{prot:C1}.
\item \textbf{Relativistic regime.} The kinetic term assumes
$v\ll c$. Relativistic generalisation replaces the Galilean kinetic
term with the Lorentzian one and reproduces the same conclusions in
the non-relativistic limit.
\end{enumerate}

\subsection{What success and failure mean}

If Protocol~\ref{prot:C1} passes, the Loschmidt operation has been
performed approximately one million times in one second of analytical
observation on the atom whose count defines the SI second, and
observed to leave the cumulative count monotone every time. The
conventional reading of Loschmidt's paradox is operationally
falsified; the framework is operationally validated.

If Protocol~\ref{prot:C1} fails, the framework is operationally
falsified. We would then need to identify which assumption breaks: the
Lagrangian (Theorem~\ref{thm:Lag}), the time-count identity
(Theorem~\ref{thm:tc}), or the monotonicity argument
(Theorem~\ref{thm:mono}).

The headline experiment is, in either outcome, decisive.

% =========================================================
\section*{Acknowledgements}
% =========================================================

This work is part of the bounded-phase-space programme at the
TU Munich Chair of Proteomics and Bioanalytics. All datasets cited
are publicly available.

% =========================================================
\bibliographystyle{plainnat}
\bibliography{references}
% =========================================================

\end{document}
